\subsection{مبانی نظری و فرمول‌بندی ریاضی}
فرض کنید ماتریس وزن اولیه یک لایه به صورت $\mathbf{W} \in \mathbb{R}^{d_{\text{out}} \times d_{\text{in}}}$ باشد. در گام نخست، تجزیه مقادیر منفرد بر روی ماتریس وزن صورت می‌گیرد:
\begin{equation}
\mathbf{W} = \mathbf{U} \mathbf{\Sigma} \mathbf{V}^\top
\end{equation}
با انتخاب رتبه هدف $r \ll \min(d_{\text{out}}, d_{\text{in}})$، تقریب کم‌مرتبه مرتبه $r$ به صورت فاکتورهای متقارن بازسازی می‌شود:
\begin{equation}
\mathbf{A} = \mathbf{U}_{:, :r} \mathbf{\Sigma}_{:r}^{1/2} \in \mathbb{R}^{d_{\text{out}} \times r}, \quad \mathbf{B} = \mathbf{\Sigma}_{:r}^{1/2} \mathbf{V}_{:, :r}^\top \in \mathbb{R}^{r \times d_{\text{in}}}
\end{equation}
ماتریس حاصل‌ضرب بازسازی‌شده کم‌رتبه برابر با $\hat{\mathbf{W}}_{\text{lr}} = \mathbf{A} \mathbf{B}$ خواهد بود. در گام دوم، نگاشت کوانتیزاسیون یکنواخت با تعداد بیت مشخص $b$ (به طور پیش‌فرض $b=8$) بر روی تانسور وزن اعمال می‌گردد. عملگر کوانتیزاسیون یکنواخت کمینه-بیشینه (\lr{Uniform Min-Max Quantization}) به صورت زیر تعریف می‌شود:
\begin{equation}
q_{\text{max}} = 2^b - 1, \quad s = \frac{\max(\mathbf{W}) - \min(\mathbf{W})}{q_{\text{max}}}
\end{equation}
\begin{equation}
\mathcal{Q}(\mathbf{W}, b) = \operatorname{clip}\left(\left\lfloor \frac{\mathbf{W} - \min(\mathbf{W})}{s} \right\rceil, 0, q_{\text{max}}\right) \cdot s + \min(\mathbf{W})
\end{equation}
سپس خطای نسبی بازسازی فروبنیوس مورد پایش قرار می‌گیرد:
\begin{equation}
\mathcal{E}_{\text{rel}} = \frac{\|\mathbf{W} - \mathcal{Q}(\hat{\mathbf{W}}_{\text{lr}}, b)\|_F}{\|\mathbf{W}\|_F + \epsilon}
\end{equation}
در صورتی که خطای نسبی بازسازی از آستانه مجاز $\mathcal{E}_{\text{max}}$ (به عنوان مثال $0.20$) فراتر رود، فرآیند تلاش مجدد تطبیقی با افزایش عمق بیتی ($b \leftarrow b+1$) و افزایش رتبه ($r \leftarrow r \times 1.25$) آغاز شده و در صورت عدم همگرایی، لایه به حالت امن اولیه بازگردانده می‌شود (\lr{Fallback}).
